\icode{E-MD-CA-L2-261114-010}
在一些长下坡路段行车道外侧时常会增设避险车道，车道表面是粗糙碎石，其作用是供下坡汽车在刹车失灵的情况下避险。如图所示为某处避险车道，现有一辆质量为\(m =\) \,\qty{2000}{kg}的汽车沿倾角\(\alpha =\) \,\qty{17}{\degree}的下坡路面行驶，当驾驶员发现刹车失灵的同时发动机失去动力，此时速度\(v_1 =\) \,\qty{10}{m/s}	，汽车继续沿下坡匀加速直行\(L =\) \,\qty{200}{m}时到达避险车道，此时速度\(v_2 =\) \,\qty{30}{m/s}（\(\sin\) \qty{17}{\degree}=0.3，\(\cos\) \qty{17}{\degree}=0.96）。求：
\begin{enumerate}
	\item 汽车下坡时的加速度大小；
	\item 汽车下坡时受到的阻力大小；
	\item 若避险车道与水平面的夹角为\qty{30}{\degree}，汽车在避险车道上受到路面的阻力是下坡公路上的 4 倍，求汽车在避险车道上运动的最大位移。
\end{enumerate}
\begin{answer}
    
\end{answer}